\section{A finite-window nonlinear inverse}
\label{nonlinear}
For vacuum unpolarized Gowdy let $V=P_\tau$ and $R=Q_\tau$.
Given $P(f-\Delta),P(f),Q(f),R(f)$, the earlier method initialized
$V$ at $f$ to the trailing secant $v=(P(f)-P(f-\Delta))/\Delta$.
Its stored values were
\[
 (v, P(f)-fv, Q(f)+R(f)/(2v_{\rm safe}), e^{2P(f)}R(f)).
\]
At a known depth the original code's clipping convention makes
those values algebraically equivalent to the observations used
here. Thus the secant restart copies $P,Q,R$ at the endpoint.
It changes the missing $V$ component. The change in backward
field history measures the performance of the restart rule.

A finite-window inverse estimates $V$ by integrating the reduced
equations back to $f-\Delta$ from a trial endpoint velocity $V_n$.
If $P_{\mathrm{trial}}(f-\Delta;V_n)$ is the resulting field, the
iteration is
\[
 V_{n+1}=V_n+
 \frac{P_{\mathrm{trial}}(f-\Delta;V_n)-P(f-\Delta)}{\Delta}.
\]
Only the specified observations are used in this iteration.
The endpoint truth velocity and earlier comparison times are
reserved for evaluation. To see why the iteration can work, write
$P_{\mathrm{trial}}(f-\Delta;V)=P(f)-\Delta V+J(V)$, where
$J(V)=\int_0^\Delta(\Delta-u)P_{\tau\tau}(f-u;V)\,du$.
For a fixed finite-dimensional discretization, a state-flow
Lipschitz bound $e^{Lu}$ and acceleration Lipschitz bound $L_F$
make $J/\Delta$ Lipschitz with constant at most
$\frac12\Delta L_F e^{L\Delta}$ on an admissible ball. An
invariant ball and a bound below one would give local convergence.
We do not use this sufficient bound as a certificate for the
simulations.

At a homogeneous polarized vacuum background,
a finite Fourier Galerkin system gives a further analytic check.
When the initial $Q$ and $Q_t$ perturbations vanish, the
linearized equation for the unmeasured $P_t$ component
is the polarized equation. If all the retained two-time
multipliers are nonzero, the derivative of the map from that
velocity to the second field snapshot is invertible.
The finite-dimensional inverse function theorem then gives
a local inverse of this reduced observation map.
The neighborhood and constants can depend on the bandwidth.

The exact true-spike family from Lim provides a local
known-answer control \cite{lim,lim2}. The corrected form for
$Q$ already occurs in Appendix A of \cite{lim2}. For fixed
time, $e^{-P(x)}$ is quadratic about the spike center.
The center can be estimated from the vertex using sampled
$P$ alone. Evaluating the two fitted quadratics there gives
the finite-window center slope and fixes the family velocity
parameter. Center values of $P$ and $Q$ then determine the
remaining family parameters. This procedure uses information
available in the original tuple and the formulas of a known
family. An off-cell center can be recovered without providing
its grid shift. The calculation is local in $x$ and does not
identify a general periodic feature as a Lim solution.

The original production example had $P=R=0$,
$V=2+1.5\cos\theta$, and $Q=\cos\theta$ at $\tau=0$.
The ordinary secant and the inverse receive identical observations
at $f-2,f$ with $f=12$ or $18$. We retain the original mask at
depth 18, while also scoring the full periodic domain.
For two states, let $d_H$ be hyperbolic distance in the $(Q,e^{-P})$
upper half-plane, $s=(V^2+e^{2P}R^2)^{1/2}$, and
$E=s^2+e^{-2\tau}(P_\theta^2+e^{2P}Q_\theta^2)$.
The score is the sampled supremum of
$\bigl(d_H^2+(\Delta s)^2+[(\Delta E)/(1+|E_{\mathrm{ref}}|)]^2\bigr)^{1/2}$.
It involves the reduced fields and their derivatives.
The comparison omits the value of the metric potential.

At $f=18$ the following scores are measured at the held-out
slice $\tau=2$. The legacy forward stencil is second order.
The two inverse columns use a fourth-order restart. The
privileged case supplies the endpoint velocity from the reference.
\begin{center}
\begin{tabular}{rccccc}
\toprule
 & \multicolumn{2}{c}{Secant, legacy stencil} &
 \multicolumn{2}{c}{Window inverse, fourth order} & Privileged \\
$N$ & Mask & Full circle & Mask & Full circle & Full circle\\
\midrule
1024 & $0.82012$ & $2.73288$ & $0.04666$ & $0.04666$ & $0.04734$\\
2048 & $0.81689$ & $3.09459$ & $0.01186$ & $0.01186$ & $0.01203$\\
4096 & $0.82215$ & $3.03402$ & $0.00298$ & $0.00298$ & $0.00302$\\
\bottomrule
\end{tabular}
\end{center}
At depth 18 the window fit residual divided by the observation
gap is below $10^{-10}$ at every tested grid. On the finest
grid the inverse gives $1.15\times10^{-3}$ and $1.55\times10^{-4}$
at $\tau=4$ and $8$. Depth 12 has no flagged component.
There the fine-grid full-circle score at $\tau=2$ falls from
$0.43207$ to $9.95\times10^{-5}$.
A same-stencil inverse is also recorded and tests inversion of its
own discretization. An analytic periodic geodesic wave-map control
provides snapshots independently of the solver. Its fourth-order
window inverse reduces the $\tau=2$, $N=1024$ score from
$0.02577$ to $1.62\times10^{-9}$.
For the two noiseless local Lim controls, the estimated half-cell
center errors were at most $1.45\times10^{-15}$ in coordinates.
Those controls test recovery within a specified family.

For a separate nonlinear check, evolve from $P=0$,
$V=2+1.4\cos\theta$, $Q=0.9\cos\theta+0.1\sin2\theta$,
and $R=0.025\cos3\theta$, with a periodic primitive of the
momentum source. Observe at $\tau=10,12$ using a fourth-order
spatial stencil. Reconstruct using the historical second-order
stencil and compare at $\tau=2$.
Apply the same reduced full-circle score.

\begin{center}
\begin{tabular}{rccc}
\toprule
$N$ & Secant & Window inverse & Exact endpoint velocity\\
\midrule
512  & $0.23690$ & $0.009585$ & $0.009876$\\
1024 & $0.23614$ & $0.002405$ & $0.002479$\\
\bottomrule
\end{tabular}
\end{center}
The matched-window score is close to that of a restart supplied
with the exact numerical endpoint velocity. On the finer grid,
adding an absolute error of $10^{-7}$ to the earlier $P$
observation in a single cosine mode gave gains of 5.17 and 8.36
in the comparison score at modes 1 and 128, respectively.
These are two directional measurements for this IVP. They
provide no infinite-band bound. The thin-feature classification
and the asymptotic regime of nonlinear unpolarized Gowdy
remain questions distinct from finite-window reconstruction.
